\chapter*{Abstract}
\addcontentsline{toc}{chapter}{Abstract}

Urban infrastructure and high-density pedestrian facilities generate continuous mechanical kinetic energy that typically dissipates unharnessed into the environment as acoustic vibration and heat. With the proliferation of distributed Internet of Things (IoT) edge nodes, wireless sensor networks, and smart transit tracking beacons, there is a compelling need for decentralized, sustainable micro-energy harvesting coupled with intelligent kinetic analytics. This report presents \textbf{StepCharge AI}, an end-to-end intelligent IoT footstep kinetic energy harvesting, storage monitoring, and machine-learning-driven biomechanical analysis system.

The hardware subsystem comprises an array of piezoelectric ceramic transducers (Lead Zirconate Titanate, PZT) embedded within a reinforced floor tile module. Dynamic mechanical compressive stress from human footfalls generates alternating high-voltage, low-current AC electrical pulses via the direct piezoelectric effect. The harvested energy is routed through a low forward-drop Schottky bridge rectifier and transient protection network into a 0.1\,F, 5.5\,V supercapacitor electrostatic storage buffer. An Espressif ESP32 microcontroller performs edge signal conditioning via 12-bit analog-to-digital converter (ADC) channels with 16$\times$ oversampling and Exponential Moving Average (EMA) filtering ($\alpha = 0.25$) on calibrated precision voltage dividers. Footstep impacts are isolated using a dual-threshold Schmitt trigger with hysteresis ($V_{\text{high}} = 0.80\,\text{V}$, $V_{\text{low}} = 0.55\,\text{V}$) and a 200--250\,ms refractory lockout window to eliminate false double-triggering.

The edge telemetry pipeline transmits encrypted time-series frames over TLS/HTTPS using SHA-256 authenticated device keys to a cloud-hosted Node.js + Express backend integrated with MongoDB Atlas. Real-time telemetry frames and footstep events are broadcast via Socket.IO room namespaces to an interactive React 19 + TypeScript web dashboard. To transform physical kinetic signals into actionable biomechanical insights, a Python FastAPI microservice extracts five time-domain kinetic features ($V_{\text{peak}}$, $V_{\text{avg}}$, $T_{\text{pulse}}$, $\Delta T_{\text{step}}$, $V_{\text{storage}}$) and classifies steps into \texttt{LIGHT}, \texttt{NORMAL}, and \texttt{HEAVY} categories using a Random Forest classifier. To prevent gait identity memorization and overly optimistic performance estimation, model evaluation enforces subject-independent Stratified Group $K$-Fold cross-validation across participants. Furthermore, a closed-loop load switching architecture ensures that demo loads (LED and low-power DC fan) are visually rendered on the dashboard only upon verified GPIO hardware readback.

\vspace{0.18in}
\noindent
\textbf{Keywords:} Energy Harvesting, Piezoelectric Transducers, Internet of Things (IoT), ESP32, Biomechanical Monitoring, Machine Learning, Random Forest, Supercapacitor Storage, Real-Time Telemetry, MongoDB Atlas.
\clearpage
